% Appendix A, Chapter 3 — robustness checks. First entries added 11 Sep 2026 (PJ: "why not get
% going on these 3.4 and 3.5" — a robustness pass on Ch3's own coding decisions). Content mirrors
% data_workspace_2026-07/ch3_pipeline_restored/scripts/67_mp_misconduct_robustness.R and
% 68_ministerial_departures_robustness.R, and the online_appendix/ch3_fig34_*.csv /
% ch3_fig35_*.csv outputs those scripts write. Regenerate/extend as further checks are run.
%
% RC17 UPDATED 14 Sep 2026: Figure 3.7's mediation model was respecified from contemporaneous to
% a sequential lag (Prof(t-2)->Trust(t-1)->Turbulence(t)) after PJ caught the mismatch between a
% mediation claim and same-year measurement. Redoing the detrending/first-differencing checks
% under the new lag structure surfaced a genuinely different result for first-differencing (a
% significant NEGATIVE relationship, not the old null) — checked for fragility (influential-year
% and HAC-lag sensitivity) before being written up, per the standing "diagnose, don't hide" rule.
% This table's caption is now referenced as \label{tab:ch03robustness} (= Table A3.1) from the
% chapter's own footnote 22, which previously had no cross-reference target.

\section{Summary of checks run}

\begin{threeparttable}
\caption{Robustness checks, Chapter Three}
\label{tab:ch03robustness}
\begin{tabular}{p{1.1cm} p{3.5cm} p{4.3cm} p{3.6cm}}
\toprule
Check & Baseline & Alternative specification & Verdict \\
\midrule
RC13 & Figure 3.4 misconduct series: Mortimore \& Blick ``Seats Forfeited'' list, $N=8$ events
(2009--2022) & Cross-checked against HC Library SN04594; two genuine events the source list missed
(Onasanya 2019, Amesbury 2025) verified independently and added &
\textbf{Corrected, not a robustness variant} --- $N$: 8 $\rightarrow$ 10; both new events fall
2009+, reinforcing rather than disrupting the ``zero events 1979--2008'' pattern \\
RC14 & Figure 3.5 reshuffle-cluster threshold: 5+ same-day appointment-endings & 3+ and 10+
thresholds & \textbf{Survives} --- standalone rate rises across 1979--2025 at all three
thresholds; only the level shifts (274 vs 216 vs 340 events, 47yr total) \\
RC15 & Figure 3.5 classification windows: 14-day lateral-move window, $\pm$3-day
government-ended window & 7-day and 30-day lateral windows; $\pm$7-day government-ended window &
\textbf{Survives} --- 47-year standalone total ranges 269--280 (baseline 274) across all four
specs; period-rate trend shape unchanged \\
RC16 & Figure 3.5 Panel C (financial scandal), $N=12$ events, no detectable break (Chow $p=0.235$,
Poisson LR $p=0.124$, both peak 1993) & Two borderline cases (Fox 2011, Siddiq 2025) added
($N=14$); plus exact small-sample tests (Poisson rate-ratio, Fisher's exact) at a pre-specified
split (2008) vs the searched peak (1993) & \textbf{Survives, and clarifies a look-elsewhere trap}
--- null holds with or without the 2 borderline cases; the pre-specified 2008 split is null by
every test ($p\geq0.46$), while the searched 1993 split looks significant \emph{only} if tested
naively without correcting for the search ($p=0.02$--$0.04$) --- see below \\
RC17 & Figure 3.7 mediation (sequential lag: Prof$_{t-2}\to$Trust$_{t-1}\to$Turbulence$_t$,
$N=45$): quadratic direct effect of professionalisation on turbulence significant ($p=0.0001$),
indirect effect via trust null ($p=0.875$) & Linear + quadratic calendar-time trend added to all
three equations; separately, a first-differenced (change-on-change) specification preserving the
same lag structure & \textbf{Mixed --- levels and change-on-change results diverge}: the direct
effect is not distinguishable from a shared secular time trend in levels ($p=0.395$
trend-adjusted, HAC $p=0.450$), matching the earlier finding; but first-differencing now finds a
significant \emph{negative} relationship instead of nothing ($p=0.034$ HAC direct effect; $p=0.050$
HAC total effect) --- opposite in sign to the positive levels relationship, and confirmed not to
be an artefact of any single influential year or the HAC lag choice (see below) \\
RC18 & Combined bad-behaviour count (Commons + Lords + local government, $N=44$ events,
1979--2025) linked to professionalisation PC1 and to the Ch2 turbulence PC1, three directions
tested: Poisson(scandal $\sim$ prof), OLS(prof $\sim$ scandal), OLS(turbulence $\sim$ scandal), all
significant at baseline ($p<0.0001$, $p<0.0001$, $p=0.007$--$0.017$ by lag respectively) &
Linear+quadratic calendar-time control and first-differencing, applied to all three directions;
quadratic/GAM non-linearity check on each & \textbf{Two directions do not survive at all; the
third (scandal $\to$ turbulence) is more real but still not causal-grade} --- professionalisation
$\leftrightarrow$ scandal is killed by both year-control ($p=0.66$--$0.81$) and first-differencing
($p=0.70$--$0.78$) in both directions, and is confirmed linear (GAM edf$=1.00$ both directions, no
curvature at all). Scandal $\to$ turbulence weakens under year-control ($p=0.19$--$0.24$, not
fully dead) but still collapses under first-differencing ($p=0.87$--$0.97$); its apparent curvature
(quadratic term $p=0.05$, GAM $p=0.02$--$0.03$ across lags) also fails to survive year-control
($p=0.13$--$0.15$) \\
\bottomrule
\end{tabular}
\begin{tablenotes}\footnotesize
\item Full detail and re-runnable scripts: \texttt{data\_workspace\_2026-07/ch3\_pipeline\_
restored/scripts/67\_mp\_misconduct\_robustness.R} (RC13), \texttt{68\_ministerial\_departures\_
robustness.R} (RC14--16), \texttt{41p\_mediation\_quadratic\_lagged.R}, \texttt{51b\_mediation\_
time\_trend\_lagged.R}, and \texttt{51c\_mediation\_firstdiff\_diagnostics.R} (RC17),
\texttt{72\_combined\_scandals\_vs\_professionalisation.R} (RC18). Outputs:
\texttt{online\_appendix/ch3\_fig34\_misconduct\_robustness.csv},
\texttt{ch3\_fig35\_threshold\_robustness.csv}, \texttt{ch3\_fig35\_window\_robustness.csv},
\texttt{ch3\_fig35\_scandal\_borderline\_robustness.csv}, \texttt{ch3\_fig35\_scandal\_exact\_
tests.csv}, \texttt{ch3\_fig36\_mediation\_trend\_robustness.csv} (superseded by
\texttt{data\_workspace\_2026-07/ch3\_pipeline\_restored/data/processed/ch3\_mediation\_time\_
trend\_results\_LAGGED.csv}).
\end{tablenotes}
\end{threeparttable}

\section{RC13 --- a source cross-check that became a correction}

Figure 3.4 (MPs losing their seat over misconduct) uses Mortimore \& Blick's ``Seats Forfeited''
list as its construct, having already rejected HC Library SN04594 (``Sitting Members imprisoned
since 1945'') as a standalone source --- SN04594 is dominated by a single clustered 1986--91
Northern Ireland Troubles-era protest-fine episode, unrelated to the personal/financial-misconduct
construct the chapter needs.

Documenting that rejection properly meant reading SN04594's own list in full, which surfaced two
events inside the 1979--2025 window that fit the chapter's construct exactly but are absent from
Mortimore \& Blick: Fiona Onasanya (2019, perverting the course of justice, removed via the
first-ever successful Recall of MPs Act 2015 petition) and Mike Amesbury (2025, common assault,
resigned under pressure after a suspended sentence). Both were verified independently of SN04594's
summary --- Onasanya against UK Parliament's own recall-petition record, Amesbury against the CPS's
own press release --- before being added.

This is a correction to the base figure, not a robustness variant: $N$ moves from 8 to 10.
Both added events fall in the 2009-or-later window, so the chapter's existing argument (misconduct
resulting in loss of seat is a recent phenomenon, zero cases 1979--2008) is reinforced rather than
disrupted. The pre-correction 8-event series is retained in
\texttt{ch3\_fig34\_misconduct\_robustness.csv} as a labelled comparison row.

\section{RC14--RC15 --- Figure 3.5's discretionary classification parameters}

The reshuffle/standalone split (Figure 3.5, Panels A--B) rests on three discretionary choices,
flagged as such when the classification was built: a same-day-endings threshold for what counts as
a coordinated reshuffle (headline 5+), a lateral-move window for excluding same-person
reappointments (14 days), and a government-ended window around each corrected handover date
($\pm$3 days). All three were tested directly rather than left as unexamined defaults.

The reshuffle threshold was already computed at 3, 5 and 10 in the classification script itself but
had not previously been surfaced as an appendix table. All three thresholds show the same rising
standalone-departure rate across 1979--2025; only the absolute level moves (216 events at the
tightest threshold, 274 at the headline 5+, 340 at the loosest), exactly as expected since a looser
threshold reclassifies some reshuffle-cluster departures as standalone.

The two window parameters were tested by rebuilding the full classification pipeline under each
alternative. Tightening the lateral-move window to 7 days or loosening it to 30 days moves the
47-year standalone total by at most 6 events either side of the 274-event baseline (280 and 269
respectively); loosening the government-ended window to $\pm$7 days moves it to 271. The rising
trend across the five sub-periods is unchanged under every variant. None of the three parameters is
load-bearing for the finding.

\section{RC16 --- the scandal panel, two borderline cases, and a look-elsewhere trap}

Figure 3.5's Panel C (financial-scandal departures) already carried an honest limitation: with only
12 events across 47 years, a mean-shift Chow test and a Poisson rare-event likelihood-ratio test
both found their peak candidate break at 1993, but neither reached significance under permutation
(Chow $p=0.235$; Poisson $p=0.124$). Two cases were deliberately left out of that 12-event series as
borderline --- Liam Fox (2011, the Werritty conflict-of-interest affair) and Tulip Siddiq (2025,
resignation over reputational association rather than a personal finding of financial impropriety).
Neither is a clean fit to Berlinski, Dewan \& Dowding's ``financial scandal'' construct, which is
specifically about a minister's \emph{own} financial impropriety.

Adding both (as a robustness variant only --- neither is added to the base figure) leaves the
substantive conclusion unchanged: still no significant break by either test (Chow $p=0.176$;
Poisson $p=0.094$), still peaking at 1993.

Because 12--14 events over 47 years is a genuinely small sample, an exact-test complement was added
alongside the search-based tests: a Poisson exact rate-ratio test and Fisher's exact test, applied
at two different splits. The first split, 2008, is the real methodological seam already present in
the data (Berlinski, Dewan \& Dowding's own audited 1979--2007 census versus this project's
opportunistic 2008--2025 search) --- a split chosen for a reason external to the data, not searched
for. It is null by every measure, with or without the two borderline cases (Poisson $p=1.00$/$0.59$;
Fisher $p=0.70$/$0.46$).

The second split, 1993, is the same year the Chow and Poisson \emph{search} procedures identified as
their best candidate. Testing an exact test naively at exactly that point --- without correcting for
having searched roughly 35 candidate years to find it --- produces an apparently significant result
(Poisson $p=0.026$; Fisher $p=0.041$; both tighten slightly further with the two borderline cases
included). This is not a second, independent confirmation of a 1993 break: it is the search
procedure's own selection effect showing up when the correction for that selection is removed. The
sup-statistic permutation tests already correct for this by permuting the entire search, which is
precisely why they --- correctly --- do not find significance at the same point. Reported here
because the discrepancy is itself instructive, not because it changes the finding: read together,
the four tests support one conclusion, that the scandal series shows no detectable structural break
at any pre-specified point, and that a search-identified candidate should never be tested again
without the correction that found it in the first place.

\section{RC17 --- Figure 3.7's mediation result under calendar-time detrending and
first-differencing}

Figure 3.7's headline finding is that professionalisation has a significant quadratic direct
relationship with turbulence (linear term $p=0.0001$, quadratic term $p=0.0007$), while the
indirect path via political trust is null (ACME $p=0.875$) --- read as evidence that
professionalisation matters, just not through trust. The model uses a sequential lag structure
(respecified 14 Sep 2026, replacing an earlier contemporaneous version that conflated a
mediation claim with same-year measurement): Professionalisation$_{t-2}\to$Political
Trust$_{t-1}\to$Turbulence$_t$, $N=45$ (1981--2025). Both professionalisation's PC1 and the
canonical turbulence composite drift upward over the 47-year span, raising the standard concern
that two independently-trending series can appear related without any real link between them,
purely because both are moving in the same direction over time.

Two independent detrending strategies were applied to check this, deliberately using different
mechanisms so that agreement between them is informative rather than an artefact of one method's
assumptions.

The first adds linear and quadratic calendar-time terms directly to all three equations. Under
this specification, professionalisation's direct relationship with turbulence in levels is no
longer significant (linear $p=0.395$, quadratic $p=0.530$; Newey--West HAC $p=0.450$/$0.584$), its
relationship with trust also disappears ($p=0.432$ OLS), and the already-null indirect effect
stays null (ACME $p=0.940$). This matches the earlier finding on the contemporaneous version of
this model: the quadratic \emph{levels} relationship does not survive controlling for a shared
secular trend.

The second strategy, first-differencing, removes any smooth trend (of whatever functional form)
by construction rather than modelling it directly, preserving the same $t-2$/$t-1$/$t$ offsets in
the differenced series. This time it does \textbf{not} reach the same conclusion: the change in
professionalisation (2 years prior) shows a significant \emph{negative} relationship with the
change in turbulence (direct effect $p=0.034$, total effect $p=0.050$, both Newey--West HAC; plain
OLS is only marginal, $p=0.099$/$0.136$) --- the opposite sign to the positive levels relationship.
This was checked for fragility before being reported, given the counter-intuitive sign flip: it is
not driven by any single influential year (dropping the three years with the highest Cook's
distance makes the result \emph{stronger}, not weaker --- OLS itself then crosses into
significance, and the HAC $p$-value falls to $0.001$), and it is stable across Newey--West lag
orders 0--4 ($p=0.029$--$0.042$ throughout), so it is not an artefact of the specific HAC lag
choice either.

A simpler, non-mediated alternative was also considered --- the plain bivariate
Professionalisation$_{t-1}\to$Turbulence$_t$ quadratic relationship (R$^2=0.458$, quadratic term
$p=5.0\times10^{-5}$, GAM edf$=2.70$) --- on the reasoning that since trust never mediates
anything in any specification, the chapter's claim could rest on this simpler relationship
instead. It does \emph{not} hold up as well as the full mediation model under the same two
checks: the trend-adjusted quadratic term is null ($p=0.25$ OLS/$0.26$ HAC, GAM $p=0.42$), and
first-differencing is also null (quadratic term $p=0.42$ OLS/$0.088$ HAC) --- unlike the full
model, there is no significant finding left over on either check. The simpler bivariate framing
was tried and set aside for this reason, not adopted.

\textbf{Verdict, and how the chapter treats it}: Figure 3.7 is retained as the unadjusted,
exploratory levels model --- the null finding on trust-mediation is unaffected by any of this, and
the quadratic direct relationship it displays in levels should still not be read as a robust or
causal effect of professionalisation on turbulence, since calendar-time detrending removes it.
But unlike the contemporaneous version of this model (and unlike the simpler bivariate
alternative above), first-differencing here does \emph{not} also support a null: it surfaces a
distinct, checked-for-robustness, negative change-on-change relationship, opposite in sign to the
positive levels association. The chapter text states the direct (levels) effect as an unadjusted
association, not a causal claim, and reports the first-differenced finding as a separate result in
its own right rather than a second confirmation of the same one.

\section{RC18 --- the combined bad-behaviour count against professionalisation and against
turbulence}

PJ, 14 Sep 2026: "is there a way to map scandals to the PCA index -- not just MPs but all of them
together (not devolved) -- need to model for zero cases." The combined series pools the Commons
misconduct series (Fig.\ 3.3, systematic Mortimore \& Blick census), the House of Lords case list,
and the local government case list (illustrative, web-compiled -- see the respective source
notes) into one annual count, $N=44$ events over 47 years, 25 of them (53\%) zero-event years.
Ministerial scandal resignations were tried and dropped from the pool (worse fit on every test,
including the direction that matters most -- see the header of
\texttt{72\_combined\_scandals\_vs\_professionalisation.R} for the full comparison). Because most
years are zero, Poisson regression is used for scandal-as-outcome models rather than OLS/
correlation, matching the change-point tests used elsewhere in this chapter.

Three directions were tested, all significant at baseline: professionalisation PC1 predicting the
combined scandal count (Poisson, IRR $\approx1.75$--$1.81$ by lag, $p<0.0001$, McFadden
$R^2=0.23$--$0.24$); the combined scandal count predicting professionalisation PC1 (OLS,
$R^2=0.37$--$0.41$, $p<0.0001$); and the combined scandal count predicting the actual Ch2
turbulence PC1 -- the book's real dependent variable, and the direction that matters most for the
chapter's argument (OLS, coefficient $0.39$--$0.44$ by lag, $p=0.007$--$0.017$,
$R^2=0.12$--$0.16$).

The same detrending discipline as RC17 was applied to all three. The professionalisation
$\leftrightarrow$ scandal relationship, in \emph{either} direction, does not survive: year-control
kills it (professionalisation $\to$ scandal $p=0.81$; scandal $\to$ professionalisation $p=0.66$),
first-differencing kills it again by an independent route ($p=0.70$ and $p=0.78$ respectively),
and a quadratic/GAM check confirms both directions are genuinely linear with zero curvature (GAM
edf$=1.00$ exactly, both directions) -- so there is nothing non-linear being missed either. This
is a textbook shared-trend artefact: both series drift upward over the period and appear related
purely because of that.

The scandal $\to$ turbulence direction is more interesting and holds up somewhat better, though
still short of a causal-grade finding. Year-control weakens it to non-significance
($p=0.19$--$0.24$ depending on lag) but does not obliterate it the way it does the
professionalisation relationship -- the coefficient direction and magnitude stay similar, they
simply lose significance with only 47 years of data and a year-trend term added. First-differencing
does fully kill it ($p=0.87$--$0.97$ at every lag), so it cannot be read as robust to a
change-on-change specification. A quadratic/GAM check found apparent curvature at every lag
(quadratic term $p=0.05$; GAM $p=0.02$--$0.03$, edf $1.8$--$2.2$), but this too fails to survive
year-control (GAM $p$ rises to $0.13$--$0.15$) and is confirmed as a trend artefact, not a real
non-linearity, the same conclusion as RC17's parallel check.

\textbf{Verdict, and how the chapter treats it}: none of the three tested directions supports a
robust, trend-independent, causal-grade relationship. The professionalisation-scandal link (both
directions) is fully explained by a shared secular trend. The scandal-turbulence link is the
closest to a real finding -- it survives year-control better than the other two links do, even
though it does not survive first-differencing -- but should be read as suggestive rather than
established. The chapter text reports the bivariate association and its curved shape honestly,
while noting (per footnote, referencing this table) that it "disappears completely when accounting
for the time dependencies of these series."
